%\input{../preamble}

%\begin{document}

In principle, we will consider more general Borel groupoids than in most usual references - namely, we will not make any assumptions about standardness or analyticity of the Borel structure at hand. See Mackey (\cite{MR0165034,MR0201562}), Ramsay (\cite{MR0281876}) and Hahn (\cite{MR496796,MR496797}).

\begin{definition}
A \emph{Borel groupoid} consists of a groupoid $\G$ endowed with a Borel (measurable) structure (i.e., a $\sigma$-algebra) such that the product and inverse maps are Borel (where $\G[0]$ and $\G[2]$ are endowed with the induced Borel structures from $\G$ and $\G\times\G$, respectively).
\end{definition}
%\begin{enumerate}
%\item[(i)] $\G[0]$ is a Borel subset of $\G$;
%\item[(i)] $\G[2]$ is a Borel subset of $\G\times\G$;
%\item[(ii)] The product and inverse maps are Borel.
%\item[(ii)] The set $\BOR(\G)$ of Borel subsets of $\G$ for which the source and range maps are Borel isomorphisms onto Borel subsets of $\G[0]$ is closed under products.
%\end{enumerate}
%We call $\BOR(\G)$ the \emph{Borel semigroup of $\G$}.
%\end{definition}

Note that we are not assuming that $\G[0]$ is Borel in $\G$, nor that $\G[2]$ is Borel in $\G\times\G$. 

\begin{example}
Let $X$ be any uncountable set, endowed with the $\sigma$-algebra of countable or co-countable subsets. Then the diagonal $\Delta_X=\left\{(x,x)\in X\times X:x\in X\right\}$ is not Borel in $X\times X$.
\begin{itemize}
\item As a unit groupoid, $X$ is a Borel groupoid, but $X^{(2)}=\Delta_X$ is not Borel in $X\times X$;
\item Let $\G=X\times X$, the coarsest equivalence relation on $X$, endowed with the product $\sigma$-algebra. Then $\G$ is a Borel groupoid, but $\G[0]=\Delta_X$ is not Borel in $\G$ (and $\G[2]$ is also not Borel in $\G\times\G$).
\end{itemize}
\end{example}

\begin{example}
Every étale groupoid with the Borel $\sigma$-algebra (generated by the topology) is a Borel groupoid.
\end{example}

\begin{example}\label{examplesigmacompactampleisborel}
If $\G$ is a $\sigma$-compact ample Hausdorff groupoid, then $\G$ is a Borel groupoid when endowed with the Baire $\sigma$-algebra (generated by compact $G_\delta$ sets, or equivalently by $\KB(\G)$ -- see remark below). Indeed, given $A\in\KB(\G)$, the inverse $A^{-1}$ also belongs to $\KB(\G)$ and so the inversion is measurable. To check that the product map $m:\G[2]\to \G$ is measurable, let $A\in\KB(\G)$. For any pair $C,D\in\KB(\G)$, the set $m^{-1}(A)\cap(C\times D)$ is a compact-open subset of $\G[2]$, and so it can be written as a finite union
\[m^{-1}(A)\cap(C\times D)=\bigcup_{i=1}^n(A_i\times B_i)\cap\G[2]\]
for some $A_i,B_i\in\KB(\G)$. In particular, $m^{-1}(A)\cap(C\times D)$ is measurable in $\G[2]$, for any $C,D\in\KB(\G)$, and $\sigma$-compactness of $\G$ allows us to conclude that $m^{-1}(A)$ is measurable in $\G[2]$.
\end{example}

\begin{remark}
If $X$ is any zero-dimensional, locally compact Hausdorff topological space and $\mathcal{B}$ is a basis for $X$ whose elements are compact-open, then the Baire $\sigma$-algebra of $X$ is generated by $\mathcal{B}$. Indeed, every element of $\mathcal{B}$ is a compact $G_\delta$, so it is a Baire set. Conversely, if $K\subseteq X$ is compact, $U\subseteq X$ is open and $K\subseteq U$, there are $A_1,\ldots,A_n\in\mathcal{B}$ such that $K\subseteq\bigcup_{i=1}^n A_i\subseteq U$. If $K$ is $G_\delta$ we conclude that $K$ is a countable intersection of finite unions of element of $\mathcal{B}$, and therefore every Baire set is in the $\sigma$-algebra generated by $\mathcal{B}$.
\end{remark}

Borel groupoids are the measurable analogues of topological groupoids, and - in the same manner as in the topological setting - we will be interested in semigroups of bisections which encode their Borel structure.

\begin{definition}\nomenclature[G]{$\BOR(\G)$}{Borel semigroup of $\G$}\nomenclature[S]{$\BOR(\G)$}{Borel semigroup of $\G$}\index{Semigroup!Borel}
Given a Borel groupoid $\G$, define $\BOR(\G)$ as the set of Borel subsets of $\G$ for which the source and range maps are Borel isomorphisms onto Borel subsets of $\G[0]$ (given the induced structure from $\G$). We call it the \emph{Borel semigroup} of $G$.
\end{definition}

Note that if $A\in\BOR(\G)$ then $A^*\in \BOR(\G)$ as well.

\begin{proposition}
$\BOR(\G)$ is closed under products: If $A,B\in\BOR(\G)$ then $AB=\left\{ab:(a,b)\in(A\times B)\cap\G[2]\right\}\in\BOR(\G)$.
\end{proposition}
\begin{proof}
If $A\in\BOR(\G)$ and $D\subseteq\G$ is Borel, then $AD$ consists of the elements of $\G$ such that when multiplied by an appropriate element of $A^*$ on the left yield an element of $D$, i.e.,
\[AD=\left\{c\in\ra^{-1}(\ra(A)):\ra|_A^{-1}(\ra(c))^{-1}c\in D\right\}\]
is Borel, because the range map $\ra$ restricts to a Borel isomorphism from $A$ to the Borel subset $\ra(A)$ of $\G[0]$, and the product and inverse maps are Borel. Similarly, $DA$ is Borel.

If $A,B\in\BOR(\G)$, then the argument above shows that $AB$ is Borel. If $D\subseteq AB$ is Borel, then
\[\ra(D)=DB^*\]
is Borel, because $B^*\in\BOR(\G)$. In particular $\ra(AB)$ is Borel in $\G[0]$ and this also proves that the range map is a Borel isomorphism from $AB$ to its image. The source map is deal with similarly, so $AB\in\BOR(\G)$.\qedhere
\end{proof}

In general, $\BOR(\G)$ is not equal to the set of bisections which are Borel subsets of $\G$, and $\G$ cannot be covered by elements of $\BOR(\G)$. Therefore we will restrain from mentioning ``Borel bisections'', except when these two notions coincide.

\begin{example}\label{exampleomega1borelisfunny}
Let $\omega_1$ be the first uncountable ordinal\ftnt{Up to order isomorphism, this is the only uncountable well-ordered set such that every bounded subset is countable. See \cite[Chapter 2]{MR1940513} and \cite[1.19]{MR0264581} for more details.}, viewed as a unit groupoid, and let $F$ be any nontrivial countable group with unit $1_F$. Let $\G=\omega_1\sqcup F$ be the coproduct groupoid. We endow $\G$ with the $\sigma$-algebra generated by sets of the forms
\[[0,\alpha]=\left\{\beta\in\omega_1:\beta\leq\alpha\right\},\qquad\alpha\in\omega_1\]
and
\[\left\{a\right\},\qquad a\in F\setminus\left\{1_F\right\}.\]
Let us check that $\G$ is a Borel groupoid: given $\alpha\in\omega_1$, $[0,\alpha]^{-1}=[0,\alpha]$ is Borel, and if $a\in F\setminus\left\{1_F\right\}$ then $\left\{a\right\}^{-1}=\left\{a^{-1}\right\}$ is also Borel, because $a^{-1}\neq 1_F$), so the inversion map is Borel.

Note, moreover, that $\G[0]=\omega_1\cup\left\{1_F\right\}=\G\setminus\bigcup_{a\in F\setminus\left\{1_F\right\}}\left\{a\right\}$ is Borel in $\G$ (and can be identified with the successor ordinal $\omega_1+1$).

Let $m:\G[2]\to \G$ be the product map. If $\alpha\in\omega_1$, then
\[m^{-1}([0,\alpha])=([0,\alpha]\times[0,\alpha])\cap\G[2]\]
is Borel in $\G[2]$, and if $a\in F\setminus\left\{1_F\right\}$, then
\[m^{-1}(\left\{a\right\})=\bigcup_{g\in F\setminus\left\{1_F,a\right\}}\left[\left(\left\{g\right\}\times\left\{g^{-1}a\right\}\right)\cup\left(\G[0]\times\left\{a\right\}\right)\cup\left(\left\{a\right\}\cup\G[0]\right)\right]\cap\G[2]\]
is also Borel in $\G[2]$. Therefore $\G$ is a Borel groupoid.

Now note that every Borel set in $\G$ either contains a set of the form
\[(\alpha,\omega_1)\cup\left\{1_F\right\}=\left\{\beta\in\omega_1:\alpha<\beta\right\}\cup\left\{1_F\right\},\qquad\alpha\in\omega_1\]
or its complement does. In particular, $\left\{1_F\right\}$ is not Borel in $\G$, and this implies that no $a\in F\setminus\left\{1_F\right\}$ is contained in any element of $\BOR(\G)$, even though $\left\{a\right\}$ is a bisection which is a Borel subset of $\G$.
\end{example}

%\begin{example}
%Let $\G$ be any groupoid, and endow it with the topology whose open sets are those $A\subseteq\G$ for which $A\cap \G[0]\in\left\{\G[0],\varnothing\right\}$. This makes $\G$ a topological groupoid (Example~\ref{exampleweirdtopologicalgroupoid}), and in fact the Borel structure is the same as the topology, so the product and inversion maps are Borel.

%However, if $\G[0]$ has at least two elements and $\G\setminus\G[0]$ is nonempty then $\G$ is not a Borel groupoid: Take any $a\in\G\setminus\G[0]$, so $\left\{a\right\}$ and $\left\{a^{-1}\right\}$ is Borel but $\mathfrak{s}(a)=\left\{\mathfrak{s}(a)\right\}$ is not.
%\end{example}

\begin{proposition}\label{propositionwheneveryborelbisectionisborelbisection}
Let $\G$ be a Borel groupoid for which there exists a countable family $\mathcal{C}\subseteq\BOR(\G)$ such that $\G\subseteq\bigcup\mathcal{C}$. Then every bisection of $\G$ which is a Borel subset belongs to $\BOR(\G)$.
\end{proposition}
\begin{proof}
If $C\in\mathcal{C}$ and $A$ is a Borel subset of $C$, then of course $A\in\BOR(\G)$, since the source and range maps on $A$ are simply the restrictions of the source and range maps on $C$ and so are Borel isomorphism onto their (Borel) images.

In general, if $A$ is a bisection of $\G$ which is a Borel subset, then $A=\bigcup_{C\in\mathcal{C}}(A\cap C)$, and each term $A\cap C$ belongs to $\BOR(\G)$, hence $A\in\BOR(\G)$.\qedhere
%\[AB=\left\{g\in CD:\ra(g)\in\ra(A),\so(g)\in\so(B)\right\}=CD\cap\ra^{-1}(\ra(A))\cap\so^{-1}(\so(B))\]
%is Borel because the source and range maps are Borel isomorphisms of $C$ and $D$ onto Borel subsets of $\G[0]$. For general $A,B\in\BOR(\G)$, we have
%\[AB=\bigcup_{C,D\in\mathcal{C}}(A\cap C)(B\cap D)\]
%which is a countable union of Borel maps and hence it is a Borel subset of $\G$. Moreover,
%\[AB=\bigcup_{E\in\mathcal{C}}\]
%and since for every $E\in\mathcal{C}$, the source map on $E$ restricts to a Borel isomorphism of $AB\cap E$ onto the Borel subset $\so(AB\cap E)$ of $\G[0]$, the source map is a Borel isomorphism of $AB$ onto the Borel subset $\so(AB)$ of $\G[0]$, The range is dealt with similarly, so we conclude that $AB\in\BOR(\G)$.
\end{proof}

%In particular
%\begin{itemize}
%\item If $\G$ is a Lindelöf étale groupoid, then $\G$ endowed with the Borel $\sigma$-algebra, generated by its open sets, is a %Borel groupoid, because the semigroup of open bisections of $\G$ covers $\G$ and thus admits a countable subcover, which %generates a countable sub-inverse semigroup.
%\item If $\G$ is a Lindelöf ample groupoid, then $\G$ endowed with the Baire $\sigma$-algebra, generated by compact $G_%\delta$ sets, or equivalently by $\KB(\G)$, is a Borel groupoid, for a similar reason as above.
%\end{itemize}
%\end{example}

%If $\G$ is an étale groupoid which admits a countable cover of open bisections (this is the case, for example, for second countable, $\sigma$-compact, or more generally Lindelöf étale groupoids), % and such that $\G[0]$ is Hausdorff,
%then $\G$ is a Borel groupoid when endowed with its Borel $\sigma$-algebra (induced by its topology).

%Indeed, $\G[0]$ is open, % (in fact clopen), and $\G[2]=\left\{(g,h)\in\G[2]:\mathfrak{s}(g)=\mathfrak{r}(h)\right\}$ is closed, so both $\G[0]$ and $\G[2]$ are
%hence Borel. The product and inverse maps are obviously Borel, so the last property is the only non-trivial one.

%Let $\left\{A_n\right\}$ be a countable collection of open bisections for which $\G=\bigcup_nA_n$. If $B,C\in\BOR(\G)$, then
%\[BC=\bigcup_{n,m}(B\cap A_n)(C\cap A_m)\]
%We know that the product map $m:\G[2]\to\G$ restricts to a homeomorphism of $(A_n\times A_m)\cap\G[2]$ onto its image, so it also restricts to a Borel isomorphism. This means that $(B\cap A_n)(C\cap A_m)=m((B\cap C)\cap(A_n\times A_m)\cap\G[2])$ is Borel for all $m,n$, and thus $BC$ is a Borel bisection of $\G$.

%We are done if we show that any Borel bisection actually belongs to $\BOR(\G)$: Let $B$ be any Borel bisection in $\G$. If $A\subseteq B$ is Borel, then $\mathfrak{s}(A)=\bigcup_n\mathfrak{s}(A\cap A_n)$ is also Borel, so the source map is a Borel isomorphism of $B$ onto its image. Note that in this case, $\mathfrak{s}^{-1}(x)=\bigcap_n A_n\cap\mathfrak{s}^{-1}(x)$ is at most countable for all $x\in\G[0]$.
%\end{example}

%Add example of étale groupoid, not covered by countably many Borel bisections. What needs to happen is that there is a Borel bisection not covered by countably many open bisections (this weaker condition is enough for the example above
%Probably

Recall that a Borel space $X$ is \emph{standard} if there is a separable (equivalently, second-countable) completely metrizable topology on $X$ which induces the initial Borel structure. Standard Borel spaces have good analytical properties and are more manageable than general Borel spaces. For a detailed discussion of standard Borel spaces (and more general classes of Borel spaces), see \cite{MR1321597}.

We recall the reader that the \emph{graph} of a function $f:X\to Y$ is $\operatorname{graph}(f)=\left\{(f(x),x):x\in X\right\}$ (see Definition \ref{definitiongraphoffunction} and the remark succeeding it).

\begin{theorem}[{\cite[Theorems 14.12 and 15.2]{MR1321597}}]\label{theoremkechrisborelmapsofstandardaregood}%{\cite[Theorem 4.5.2]{MR1619545}}
Let $X,Y$ be standard Borel spaces and $f:X\to Y$. Then $f$ is Borel if and only if $\operatorname{graph}(f)$ is Borel in $Y\times X$, and in this case, if $A\subseteq X$ is Borel and $f|_A$ is injective then $f(A)$ is Borel.
\end{theorem}

Let $\pi_X$ and $\pi_Y$ be the projection maps of $X\times Y$:
\[\pi_X(x,y)=x,\qquad\pi_Y(x,y)=y.\]
A subset $P\subseteq X\times Y$ is called a \emph{graph} if $\pi_Y$ is injective on $P$.% By the theorem above, Borel graphs are in one-to-one correspondence to Borel functions $f:\dom(f)\to X$, where $\dom(f)\subseteq Y$ is Borel. Namely, to such a function $f$ we associate its graph $\operatorname{graph}(f)=\left\{(x,f(x)):x\in\dom(f)\right\}$.
%
One important result in the study of standard Borel spaces is the \emph{Lusin-Novikov (Uniformization) Theorem}.

\begin{theorem}[{\cite[Theorem 18.10]{MR1321597}}]\index{Theorem!Lusin-Novikov}\label{theoremlusinnovikov}%{\cite[Theorem 5.10.3]{MR1619545}}
Let $X,Y$ be standard Borel spaces and $P\subseteq X\times Y$ a Borel subset such that for all $x\in X$, the \emph{section}
\[P_y=\left\{x\in X:(x,y)\in P\right\}\]
is (at most) countable. Then $P$ is a countable union of Borel graphs.
\end{theorem}

We will be interested in the following equivalent form of the Lusin-Novikov Theorem (compare with Exercises 18.14 and 18.15 of \cite{MR1321597}):

\begin{theorem}\label{theoremalteredlusinnovikov}\index{Theorem!Lusin-Novikov}\index{Countable-to-one function}
Let $X$ and $Y$ be two standard Borel spaces and $f:X\to Y$ a \emph{countable-to-one} function, that is, $f^{-1}(y)$ is (at most) countable for all $y\in Y$. Then there is a Borel partition $\left\{X_n\right\}_n$ of $X$ such that the restriction $f|_{X_n}:X_n\to f(X_n)$ is a Borel isomorphism -- that is $f|_{X_n}$ is bijective with Borel inverse -- for all $n$ (and $f(X_n)$ is Borel in $Y$ for all $n$).

In particular if $A\subseteq X$ is Borel then $f(A)$ is Borel in $Y$.
\end{theorem}
\begin{proof}
Let $P=\left\{(x,f(x))\in X\times Y:x\in X\right\}$, which is Borel since it is the opposite of $\operatorname{graph}(f)$, and all sections $P_y$ ($y\in Y$) are countable (in fact, $P_y=f^{-1}(y)\times\left\{y\right\}$). By the Lusin-Novikov Theorem, there are Borel graphs $P_n\subseteq P$ such that $P=\bigcup_{n=1}^\infty P_n$, and of course we can assume the $P_n$ are pairwise disjoint, by substituting $P_n$ by $P_n\setminus\bigcup_{i=1}^{n-1}P_i$ for $n\geq 2$, if necessary.

The projection $\pi_X:P\to X$ is a bijection, so we can take $X_n=\pi_X(P_n)$. Theorem \ref{theoremkechrisborelmapsofstandardaregood} implies both that each of the sets $X_n$ and $f(X_n)=\pi_Y(P_n)$ are Borel in $X$ and $Y$, respectively, and that $f|_{X_n}$ is a Borel isomorphism.\qedhere
\end{proof}

\begin{remark}
The usual form of the Lusin-Novikov Theorem follows from the one above, since if $P$ has countable sections $P_y$ ($y\in Y$) then the projection $\pi_Y$ is countable-to-one on $P$, and the partition $\left\{P_n\right\}_n$ yielded by Theorem \ref{theoremalteredlusinnovikov} will be a partition by Borel graphs.
\end{remark}

\begin{corollary}\label{corollarysurjectiveborelmapshavesections}
If $X$ and $Y$ are standard Borel spaces and $f:X\to Y$ is a countable-to-one surjective Borel map, then there is a Borel map $g:Y\to X$ such that $fg=\id_Y$, that is, $f$ admits a \emph{Borel section} (equivalently, there is a Borel subset $X'$ of $X$ such that the restriction $f|_{X'}$ is a Borel isomorphism onto $Y$).
\end{corollary}
\begin{proof}
Let $\left\{X_n\right\}_n$ be a countable Borel partition of $X$ such that $f|_{X_n}$ is a Borel isomorphism onto its image for all $n$. In particular, $Y=\bigcup_n f(X_n)$, so we define $g=f|_{X_1}^{-1}$ on $f(X_1)$, and $g=f|_{X_n}^{-1}$ on $X_n\setminus\bigcup_{i=1}^{n-1} f(X_i)$ for all $n\geq 2$.\qedhere
\end{proof}

%The most problematic property of Borel groupoids is the closure of Borel bisections under products, but when we consider only standard Borel structures and a countability condition this is no longer a problem.

\begin{definition}[\cite{MR3142029,MR1799683}]\index{Groupoid!$\ra$-discrete}
A groupoid $\G$ is \emph{$\ra$-discrete} if $\G^x=\ra^{-1}(x)$ is (at most) countable for all $x\in\G[0]$ (equivalently, $\G_x=\so^{-1}(x)$ is countable for all $x\in\G[0]$).
\end{definition}

\begin{remark}
In \cite[Definition 2.6]{MR584266}, Renault uses $\ra$-discreteness in the topological setting to mean that the unit space $\G$ of a locally compact Hausdorff groupoid is open. If $\G$ is second-countable and $\G[0]$ is open then $\G_x$ is countable for all $x\in X$, but the converse is not true even in the case of a second-countable groupoid (see Example \ref{exampledifferenttopologyonZ}). The definition we adopt is of more common usage in the measurable setting (such as in \cite{MR1799683}).
\end{remark}

An immediate consequence of the Lusin-Novikov Theorem (\ref{theoremalteredlusinnovikov}) and Proposition \ref{propositionwheneveryborelbisectionisborelbisection} is the following:

\begin{proposition}\label{radiscretestandardiscoveredbycountablymanyborg}
If $\G$ is a standard $\ra$-discrete Borel groupoid, then $\G$ is covered by countably many elements of $\BOR(\G)$, and every Borel subset of $\G$ which is a bisection belongs to $\BOR(\G)$.%$Suppose $\G$ is an $\ra$-discrete groupoid endowed with a standard Borel structure for which the product and inversion maps are Borel, then $\G$ is a Borel groupoid.
\end{proposition}
%\begin{proof}
%First, note that the source map $\so:a\mapsto a^{-1}a$ is Borel, so $\so(\G)=\G[0]$ is Borel in $\G$ and in fact a standard Borel space. The map $(\ra,\so):\G\to\G[0]\times\G[0]$ is countable-to-one, so by the Lusin-Novikov Theorem \ref{theoremalteredlusinnovikov} we can find a Borel partition of $\G$ by subsets $A_n\subseteq\G$ such that the source and range maps are Borel isomorphisms of $A_n$ onto their images, that is, $A_n\in\BOR(\G)$.
%
%The proof that every $\BOR(\G)$ coincides with the set of all Borel bisections is similar to that in Example~\ref{exampleetalegroupoidsborel}. Moreover, $\so,\ra:\G\to\G[0]$ are Borel maps, and $\G[2]=\left\{(a,b)\in\G[2]:\so(a)=\ra(b)\right\}$ is also Borel in $\G[0]\times\G[0]$ (because $\G[0]$ is standard) and the product map $m:\G[2]\to\G$ is countable-to-one, since $m^{-1}(a)\subseteq\G^{\ra(a)}\times\G_{\so(a)}$, and in particular it maps Borel sets to Borel sets. So if $B,C\in\BOR(\G)$ then $BC=m((B\times C)\cap\G[2])$ is also a Borel bisection.\qedhere
%\end{proof}

\begin{remark}
If $X$ is any (possibly non-standard) Borel space, seen as a Borel unit groupoid, then $\BOR(\G)$ is simply the family of Borel subsets of $X$, and $X$ is covered by $\left\{X\right\}$. Thus standardness of a Borel groupoid is not in general a consequence of $\G$ being covered by countably many elements of $\BOR(\G)$.
\end{remark}

\begin{example}\label{exampleborelstructureequivalencerelationisunique}
If $X$ is a standard Borel space and $R$ is a $\ra$-discrete equivalence relation on $X$, then the only standard groupoid Borel structure on $R$ which induces the initial Borel structure of $X=R^{(0)}$ is the one coming from the product $X\times X$. This follows from the Lusin-Novikov Theorem applied to the inclusion map $R\hookrightarrow X\times X$.

In this case, $R$ is a Borel subset of $X\times X$, and under the usual identification of bisections of $R$ as partial maps on $X$ preserving $R$-equivalence classes (Example \ref{examplebisectionsofequivalencerelationarefunctions}), elements of $\BOR(R)$ corresponds partial Borel automorphisms $f:A\to B$, where $A,B\subseteq X$ are Borel.
\end{example}

\begin{example}
Every countable group is standard and $\ra$-discrete as a groupoid, with the discrete $\sigma$-algebra..
\end{example}

\begin{example}
More generally, if $G$ is a countable group acting via Borel automorphisms on a standard Borel space $X$, then the transformation groupoid $G\ltimes X$ and the orbit groupoid $\mathcal{R}(G\ltimes X)$ are standard and $\ra$-discrete (where we endow $G$ with the discrete $\sigma$-algebra).
\end{example}

\begin{theorem}\label{theoremquotientfrommeasuredgroupoidtoorbitrelation}
Let $\G$ be an $\ra$-discrete standard Borel groupoid. Let $R=\mathcal{R}(\G)=(\ra,\so)(\G)=\left\{(\ra(a),\so(a)):a\in\G\right\}$ be the orbit equivalence relation on $\G[0]$, endowed with the $\sigma$-algebra coming from the product $\sigma$-algebra of $\G[0]\times\G[0]$. Then $R$ is also $\ra$-discrete and standard Borel, hence an $\ra$-discrete Borel groupoid on its own right. Moreover, there is a map $\rho:\BOR(R)\to\BOR(\G)$ such that $(\ra,\so)\circ\rho=\id_{\BOR(R)}$.
\end{theorem}
\begin{proof}
As usual, we identify $\G[0]$ and $R^{(0)}$. If we denote by $\so_{\G}$ and $\so_R$ the source maps and on $\G$ and $R$, respectively, and similarly for the range maps, then for all $x\in \G[0]$,
\[\ra_R^{-1}(x)=(\ra_{\G},\so_{\G})\left(\left(\ra_{\G}\right)^{-1}(x)\right)\]
so $R$ is $\ra$-discrete. The structure maps of $R$ are all Borel, and $R$ is a Borel image of the standard Borel space $\G$ so $R$ is standard as well.% and hence discrete.

Since $(\ra_{\G},\so_{\G}):\G\to R$ is surjective, let $g:R\to\G$ be any Borel section, by Corollary \ref{corollarysurjectiveborelmapshavesections}, so that $g$ is injective and induces, by Theorem \ref{theoremkechrisborelmapsofstandardaregood}, a map $\rho:\BOR(R)\to\BOR(\G)$, $\rho(A)=g(A)$.\qedhere
\end{proof}

The following is also an important result, which proves that every principal standard discrete groupoid is actually the orbit groupoid of some Borel group action.

\begin{corollary}[{Feldman-Moore, \cite{MR0578656}}]\label{corollaryfeldmanmoore}
If $X$ is a standard Borel space and $R\subseteq X\times X$ is an $\ra$-discrete Borel equivalence relation (that is $R$ is Borel in $X\times X$), then there is a countable group $G$ and an action of $G$ on $X$ by Borel automorphisms such that $R=\left\{(gx,x):g\in G\right\}$
\end{corollary}
\begin{proof}
Since $X$ is standard, we endow it with a completely metrizable separable, and so second-countable and Hausdorff, topology which induces the Borel structure, and so the set $\left\{(x,y)\in X\times X:x\neq y\right\}$ is open and can be written as a countable union of open rectangles sets $E_n=A_n\times B_n$, where $A_n,B_n\subseteq X$ are open and $A_n\cap B_n=\varnothing$.

By the Lusin-Novikov Theorem \ref{theoremlusinnovikov}, we can partition $R$ into subsets $R_n$ for which the source and range maps $\so,\ra:R_n\to X$, $\so(y,x)=x$, $\ra(y,x)=y$, are Borel isomorphisms onto their images. For every pair of positive integers $m,n$, set
\[\gamma_{m,n}(x)=\begin{cases}
\ra(\so|_{R_n}^{-1}(x)),&\text{if }x\in A_m\cap\so(R_n),\\
\so(\ra|_{R_n}^{-1}(x)),&\text{if }x\in B_m\cap\ra(R_n),\\
x,&\text{otherwise.}
\end{cases}\]
In other words, $\gamma_{m,n}(x)=y$ if $(x,y)$ belongs to $R_n\cap((A_m\times B_m)\cup (B_m\times A_m))$ and $\gamma(x)=x$ if there is no such $y$. Then $\gamma_{m,n}$ is a Borel automorphism of $X$. Letting $G$ be the group generated by $\gamma_{m,n}$, endowed with the canonical action on $X$ ($gx=g(x)$) we obtain the desired result.\qedhere
\end{proof}

\begin{remark}
Since the generators $\gamma_{m,n}$ of the group above satisfy $\gamma_{m,n}^2=\id$, $G$ is a homomorphic image of a countable free product of the group $\mathbb{Z}/(2\mathbb{Z})$ of order $2$.
\end{remark}

\begin{definition}\label{definitionpmpgroupoid}\index{Measure!invariant}\index{Groupoid!Probability measure-preserving}
Let $\G$ be a Borel groupoid and $\mu$ a Borel measure on $\G[0]$. We say that $\mu$ is \emph{invariant} if $\mu(\mathfrak{s}(A))=\mu(\mathfrak{r}(A))$ for all $A\in\BOR(\G)$. A \emph{probability measure-preserving} groupoid is a Borel groupoid endowed with an invariant probability measure $\mu$ on $\G[0]$, which we will sometimes denote as $(\G,\mu)$ whenever $\mu$ needs to be written explicitly.
\end{definition}

\begin{remark}
If $\G$ is covered by countably many elements of $\BOR(\G)$, then $\mu$ is an invariant measure if and only if
\[\int_{\G[0]}\#\left(A\cap\G^x\right)d\mu(x)=\int_{\G[0]}\#\left(A\cap\G_x\right)d\mu(x)\]
for every Borel $A\subseteq\G[0]$. The proof that these integrals are well-defined follows the same lines as in the proof of \cite[Theorem 2]{MR0578656}. However, we will only need the invariance described in the definition above.
\end{remark}

We recall from Propositions \ref{propositionwheneveryborelbisectionisborelbisection} and \ref{radiscretestandardiscoveredbycountablymanyborg} that if $\G$ is a $\ra$-discrete standard Borel groupoid then $\BOR(\G)$ coincides with the set of Borel subsets of $\G$ which are also bisections.

\begin{proposition}\label{propositionmeasureinvariancefororbitrelation}
If $\G$ is a $\ra$-discrete standard Borel groupoid and $\mu$ is a measure on $\G[0]$, then $\mu$ is invariant for $\G$ if and only if $\mu$ is invariant for the orbit relation $\mathcal{R}(\G)$.
\end{proposition}
\begin{proof}
By Proposition \ref{radiscretestandardiscoveredbycountablymanyborg},
\[\BOR(\mathcal{R}(\G))=(\ra,\so)(\BOR(\G))\]
and if $A\in\BOR(\G)$, and $F=(\ra,\so)(A)$, then
\[\so(A)=\so(F)\qquad\text{and}\qquad\ra(A)=\ra(F)\]
from which the desired result readily follows.\qedhere
\end{proof}

For the case of group actions, see Proposition \ref{propositionactionsmeasurepreserving}

Let $(\G,\mu)$ be a probability measure-preserving groupoid. We define the \emph{null ideal} of $(\G,\mu)$ as
\[\mathbf{Null}(\G,\mu)=\left\{A\in\BOR(\G):\mu(\mathfrak{s}(\theta))=0\right\}.\]
Note that $\mathbf{Null}(\G,\mu)$ is a $\lor$-ideal of $\BOR(\G)$.

\begin{definition}\nomenclature[G]{$\MEAS(\G),\MEAS(\G,\mu)$}{Measured semigroup of a $(\G,\mu)$}\nomenclature[S]{$\MEAS(\G),\MEAS(\G,\mu)$}{Measured semigroup of a $(\G,\mu)$}\index{Semigroup!measurable}\label{definitionmeasuredsemigroup}
The \emph{measured semigroup} of a probability measure-preserving groupoid $(\G,\mu)$ is the semigroup quotient
\[\MEAS(\G,\mu)=\BOR(\G)/\mathbf{Null}(\G,\mu)\]
which we simply denote by $\MEAS(\G)$ when the measure $\mu$ is implicit and there is no risk of confusion.
\end{definition}

By Theorem \ref{theoremquotientofbooleanisboolean}, $\MEAS(\G,\mu)$ is a Boolean inverse semigroup, and in Section \ref{sectionncloomissikorski} we will see that it is actually $\sigma$-complete.

When there is no risk of confusion, we will not usually make a distinction between $\BOR(\G)$ and $\mathbf{Meas}(\G,\mu)$, in the same manner that one usually sees elements of an $L^1$ space on a probability space $(X,\mu)$ as functions on $X$, instead of classes of functions which agree $\mu$-a.e.

\begin{example}\label{exampleinterpretationofmeasr}
If $R$ is a $\ra$-discrete, probability measure-preserving Borel equivalence relation on a a standard probability space $(X,\mu)$ (see Example \ref{exampleborelstructureequivalencerelationisunique}) then $\MEAS(R)$ is isomorphic to the semigroup of partial Borel automorphisms of $X$ preserving $R$-equivalence classes, modulo identifying functions which are defined and coincide a.e.\ on the union of their domains.
\end{example}


%\end{document}